\documentclass[10pt,a4paper]{amsart}
\usepackage{amsmath,amssymb,amsthm,graphicx}
% Retained preamble for the Stacks Project chapter divisors.tex (Divisors).
% The chapter file pinned-original-source.tex carries only \input{preamble} on its first line, so no
% preamble text lives inside the chapter itself.  The Stacks Project uses one single shared preamble for
% every chapter; the snapshot stages it verbatim in sources/stacks-chow/preamble.tex.  The block below is
% copied mechanically, line for line, from that file's lines 14-18 (the Xy-pic setup) and lines 154-217
% (the native theorem declarations together with the mathematical macros).  Nothing was added, removed or
% redefined.  The wrapper class and the amsmath/amssymb/amsthm/graphicx base packages are supplied by the
% assembler, so the class/\externaldocument part of the shared preamble is not repeated here.
% For commutative diagrams we use Xy-pic
\usepackage[all]{xy}

% We use 2cell for 2-commutative diagrams.
\xyoption{2cell}
% Theorem environments.
%
\theoremstyle{plain}
\newtheorem{theorem}[subsection]{Theorem}
\newtheorem{proposition}[subsection]{Proposition}
\newtheorem{lemma}[subsection]{Lemma}

\theoremstyle{definition}
\newtheorem{definition}[subsection]{Definition}
\newtheorem{example}[subsection]{Example}
\newtheorem{exercise}[subsection]{Exercise}
\newtheorem{situation}[subsection]{Situation}

\theoremstyle{remark}
\newtheorem{remark}[subsection]{Remark}
\newtheorem{remarks}[subsection]{Remarks}

\numberwithin{equation}{subsection}

% Macros
%
\def\lim{\mathop{\mathrm{lim}}\nolimits}
\def\colim{\mathop{\mathrm{colim}}\nolimits}
\def\Spec{\mathop{\mathrm{Spec}}}
\def\Hom{\mathop{\mathrm{Hom}}\nolimits}
\def\Ext{\mathop{\mathrm{Ext}}\nolimits}
\def\SheafHom{\mathop{\mathcal{H}\!\mathit{om}}\nolimits}
\def\SheafExt{\mathop{\mathcal{E}\!\mathit{xt}}\nolimits}
\def\Sch{\mathit{Sch}}
\def\Mor{\mathop{\mathrm{Mor}}\nolimits}
\def\Ob{\mathop{\mathrm{Ob}}\nolimits}
\def\Sh{\mathop{\mathit{Sh}}\nolimits}
\def\NL{\mathop{N\!L}\nolimits}
\def\CH{\mathop{\mathrm{CH}}\nolimits}
\def\proetale{{pro\text{-}\acute{e}tale}}
\def\etale{{\acute{e}tale}}
\def\QCoh{\mathit{QCoh}}
\def\Ker{\mathop{\mathrm{Ker}}}
\def\Im{\mathop{\mathrm{Im}}}
\def\Coker{\mathop{\mathrm{Coker}}}
\def\Coim{\mathop{\mathrm{Coim}}}

% Boxtimes
%
\DeclareMathSymbol{\boxtimes}{\mathbin}{AMSa}{"02}

%
% Macros for moduli stacks/spaces
%
\def\QCohstack{\mathcal{QC}\!\mathit{oh}}
\def\Cohstack{\mathcal{C}\!\mathit{oh}}
\def\Spacesstack{\mathcal{S}\!\mathit{paces}}
\def\Quotfunctor{\mathrm{Quot}}
\def\Hilbfunctor{\mathrm{Hilb}}
\def\Curvesstack{\mathcal{C}\!\mathit{urves}}
\def\Polarizedstack{\mathcal{P}\!\mathit{olarized}}
\def\Complexesstack{\mathcal{C}\!\mathit{omplexes}}
% \Pic is the operator that assigns to X its picard group, usage \Pic(X)
% \Picardstack_{X/B} denotes the Picard stack of X over B
% \Picardfunctor_{X/B} denotes the Picard functor of X over B
\def\Pic{\mathop{\mathrm{Pic}}\nolimits}
\def\Picardstack{\mathcal{P}\!\mathit{ic}}
\def\Picardfunctor{\mathrm{Pic}}
\def\Deformationcategory{\mathcal{D}\!\mathit{ef}}

% --- appended 2026-10-12 (Chief): the project own definitions of the reference/slogan/history
% display environments, copied mechanically from the shared Stacks preamble lines 8-12
% (\usepackage{verbatim} plus the three \newenvironment lines, which the chapter preamble keeps
% verbatim from the project preamble). Nothing added, removed or redefined: these environments
% carry no mathematical content and the project itself typesets them as comments. Without these
% three lines a counted statement whose own text uses them fails with Environment slogan undefined.
\usepackage{verbatim}
\newenvironment{reference}{\comment}{\endcomment}
\newenvironment{slogan}{\comment}{\endcomment}
\newenvironment{history}{\comment}{\endcomment}
\usepackage{cleveref}
\crefname{originallemma}{original lemma}{original lemmas}
\Crefname{originallemma}{Original lemma}{Original lemmas}
\makeatletter
\newlabel{algebra-lemma-order-vanishing-determinant}{{algebra-lemma-order-vanishing-determinant}{source}{}{sourcepointer-1}{}}
\newlabel{algebra-lemma-order-vanishing-determinant@cref}{{[originallemma][0][]algebra-lemma-order-vanishing-determinant}{[1][0][]source}{}{}{}}
\makeatother
\crefname{originallemma}{original lemma}{original lemmas}
\Crefname{originallemma}{Original lemma}{Original lemmas}
\makeatletter
\newlabel{lemma-reflexive-over-regular-dim-2}{{0B3N}{source}{}{sourcepointer-2}{}}
\newlabel{lemma-reflexive-over-regular-dim-2@cref}{{[originallemma][0][]0B3N}{[1][0][]source}{}{}{}}
\makeatother
\crefname{originallemma}{original lemma}{original lemmas}
\Crefname{originallemma}{Original lemma}{Original lemmas}
\makeatletter
\newlabel{more-algebra-lemma-generalized-valuation-ring-modules}{{more-algebra-lemma-generalized-valuation-ring-modules}{source}{}{sourcepointer-3}{}}
\newlabel{more-algebra-lemma-generalized-valuation-ring-modules@cref}{{[originallemma][0][]more-algebra-lemma-generalized-valuation-ring-modules}{[1][0][]source}{}{}{}}
\makeatother
\crefname{originallemma}{original lemma}{original lemmas}
\Crefname{originallemma}{Original lemma}{Original lemmas}
\makeatletter
\newlabel{more-algebra-lemma-modules-PID}{{more-algebra-lemma-modules-PID}{source}{}{sourcepointer-4}{}}
\newlabel{more-algebra-lemma-modules-PID@cref}{{[originallemma][0][]more-algebra-lemma-modules-PID}{[1][0][]source}{}{}{}}
\makeatother
\title{Along an effective Cartier divisor on a regular surface the determinant of an exact sequence of rank r sheaves changes by a bounded twist}
\author{The Stacks Project Authors}
\date{stacks-project:divisors.tex:local-snapshot-20261012; GNU FDL 1.2 or later}
\begin{document}
\maketitle
\noindent\textit{Editorial note.} All delivered mathematical text is the author's own, copied from the pinned Stacks Project chapter file pinned-original-source.tex; no statement or proof step was rewritten, reordered or removed.\par\medskip
\noindent\textit{Editorial note.} Following the multi-claim convention for this source, source\_url is the source-level URL https://stacks.math.columbia.edu/ and the source label lemma-wedge-product-ses is carried in claim\_key and here rather than in a per-claim tag URL.\par\medskip
\noindent\textit{Editorial note.} The preamble is the chapter's shared Stacks Project preamble, retained as sources/stacks-divisors/retained-preamble.tex: a line-for-line excerpt of the shared Stacks preamble staged with this snapshot. No macro was invented or redefined, and this card adds no preamble text of its own.\par\medskip
\noindent\textit{Editorial note.} Statement span 3276-3295 and proof span 3297-3331 come from the halfA pool card operations/pool-cards/stacks-divisors-halfA.json and were re-read in the pinned chapter file before this card was written.\par\medskip
\noindent\textit{Editorial note.} The proof contains a displayed choice between two cited structure theorems for modules over a discrete valuation ring (an 'or' between two labels); the author's wording and both labels are preserved, and both are handled as external navigation pointers.\par\medskip
\noindent\textit{Editorial note.} 4 cross-references occur inside the counted spans; they are listed in external\_references as navigation pointers only.\par\medskip
\noindent\textit{Allowed standard prerequisites (editorial).}\begin{itemize}
\item Reflexive sheaves over a regular scheme of dimension two on a nonempty open where a coherent sheaf is locally free of rank r have invertible determinant, cited from the same chapter. Cited at proof line 3298. Makes the two determinants invertible, so that the canonical map between them can be compared as a section of an invertible sheaf.
\item The structure theorem for finite modules over a discrete valuation ring (or over a PID), cited from the more-algebra chapter, and the fact that dividing a finite module by its torsion does not change its double dual. Cited at proof lines 3312-3315 and used at 3316-3322. Reduces the algebraic problem to free modules of rank r, where the determinant is available.
\item The vanishing order of the determinant of a map between free modules equals the length of the cokernel, cited from the algebra chapter. Cited at proof lines 3329-3331. Gives the exact vanishing order s of the determinant in the reduced case, which is at most the minimum of s and r.
\end{itemize}\medskip
\ifdefined\hypertarget\hypertarget{sourcepointer-1}{}\fi
\noindent\textit{Source reference.} Original lemma algebra-lemma-order-vanishing-determinant: sources/stacks-divisors/pinned-original-source.tex, source lines 3330--3330. Same-chapter reference used in the retained author text at source line(s) 3330. It is supplied only as a navigation pointer so that the retained text compiles; the counted item delivers the author's own statement and the author's complete proof, and no part of the referenced item is copied in or substituted for the counted argument.\par\medskip
\ifdefined\hypertarget\hypertarget{sourcepointer-2}{}\fi
\noindent\textit{Source reference.} Original lemma 0B3N: focused-seeds/AGP/sources/stacks-snapshot-20261011/tags.txt, source lines 13971--13971. Same-chapter reference used in the retained author text at source line(s) 3298. It is supplied only as a navigation pointer so that the retained text compiles; the counted item delivers the author's own statement and the author's complete proof, and no part of the referenced item is copied in or substituted for the counted argument.\par\medskip
\ifdefined\hypertarget\hypertarget{sourcepointer-3}{}\fi
\noindent\textit{Source reference.} Original lemma more-algebra-lemma-generalized-valuation-ring-modules: sources/stacks-divisors/pinned-original-source.tex, source lines 3319--3319. Same-chapter reference used in the retained author text at source line(s) 3319. It is supplied only as a navigation pointer so that the retained text compiles; the counted item delivers the author's own statement and the author's complete proof, and no part of the referenced item is copied in or substituted for the counted argument.\par\medskip
\ifdefined\hypertarget\hypertarget{sourcepointer-4}{}\fi
\noindent\textit{Source reference.} Original lemma more-algebra-lemma-modules-PID: sources/stacks-divisors/pinned-original-source.tex, source lines 3320--3320. Same-chapter reference used in the retained author text at source line(s) 3320. It is supplied only as a navigation pointer so that the retained text compiles; the counted item delivers the author's own statement and the author's complete proof, and no part of the referenced item is copied in or substituted for the counted argument.\par\medskip
\begin{lemma}
\label{lemma-wedge-product-ses}
Let $X$ be an integral regular scheme of dimension $2$.
Let $i : D \to X$ be the immersion of an effective Cartier divisor.
Let $\mathcal{F} \to \mathcal{F}' \to i_*\mathcal{G} \to 0$
be an exact sequence of coherent $\mathcal{O}_X$-modules.
Assume
\begin{enumerate}
\item $\mathcal{F}, \mathcal{F}'$ are locally free of rank $r$ on a nonempty
open of $X$,
\item $D$ is an integral scheme,
\item $\mathcal{G}$ is a finite locally free $\mathcal{O}_D$-module
of rank $s$.
\end{enumerate}
Then $\mathcal{L} = (\wedge^r\mathcal{F})^{**}$ and
$\mathcal{L}' = (\wedge^r \mathcal{F}')^{**}$
are invertible $\mathcal{O}_X$-modules and
$\mathcal{L}' \cong \mathcal{L}(k D)$ for some
$k \in \{0, \ldots, \min(s, r)\}$.
\end{lemma}

\begin{proof}
The first statement follows from Lemma \ref{lemma-reflexive-over-regular-dim-2}
as assumption (1) implies that $\mathcal{L}$ and $\mathcal{L}'$
have rank $1$. Taking $\wedge^r$ and double duals are functors, hence
we obtain a canonical map $\sigma : \mathcal{L} \to \mathcal{L}'$
which is an isomorphism over the nonempty open of (1), hence
nonzero. To finish the proof, it suffices to see that
$\sigma$ viewed as a global section of
$\mathcal{L}' \otimes \mathcal{L}^{\otimes -1}$ does not
vanish at any codimension point of $X$, except at the generic
point of $D$ and there with vanishing order at most $\min(s, r)$.

\medskip\noindent
Translated into algebra, we arrive at the following problem:
Let $(A, \mathfrak m, \kappa)$ be a discrete valuation ring
with fraction field $K$. Let $M \to M' \to N \to 0$ be an exact sequence
of finite $A$-modules with $\dim_K(M \otimes K) = \dim_K(M' \otimes K) = r$
and with $N \cong \kappa^{\oplus s}$. Show that the induced map
$L = \wedge^r(M)^{**} \to L' = \wedge^r(M')^{**}$ vanishes to
order at most $\min(s, r)$. We will use the structure theorem for
modules over $A$, see
More on Algebra, Lemma
\ref{more-algebra-lemma-generalized-valuation-ring-modules} or
\ref{more-algebra-lemma-modules-PID}.
Dividing out a finite $A$-module by a torsion submodule does not
change the double dual.
Thus we may replace $M$ by $M/M_{tors}$ and $M'$ by
$M'/\Im(M_{tors} \to M')$ and assume that $M$ is torsion free.
Then $M \to M'$ is injective and $M'_{tors} \to N$ is injective.
Hence we may replace $M'$ by $M'/M'_{tors}$ and $N$ by $N/M'_{tors}$.
Thus we reduce to the case where $M$ and $M'$ are free of rank $r$
and $N \cong \kappa^{\oplus s}$. In this case $\sigma$
is the determinant of $M \to M'$ and vanishes to order $s$
for example by Algebra, Lemma \ref{algebra-lemma-order-vanishing-determinant}.
\end{proof}

\end{document}
